% !TeX root = ../../Book2.tex
\section{Riesz–Thorin 插值}
\label{b2:sec:1704}

线性算子可沿复参数改变输入和输出的指数，从而得到具有精确几何平均常数的插值估计。为使用这项方法，需要复变函数的最大模原理；本节把所需的条带估计单独证明。

\subsection{解析族}
\label{b2:subsec:170401}

下面采用\term[hadamard-sanxian-yinli]{Hadamard 三线引理}[Hadamard three-lines lemma]的下述版本。

\begin{lemma}[Hadamard 三线估计]\label{b2:thm:17-three-lines}
设 \(F\) 在 \(0<\operatorname{Re}z<1\) 全纯，在闭条带连续且有界，边界满足 \( |F(\mathrm it)|\leq M_0\)、\( |F(1+it)|\leq M_1\)。则 \(0<\theta<1\) 时
\[
 |F(\theta)|\leq M_0^{1-\theta}M_1^\theta.
\]
\end{lemma}
\begin{proof}
先设 \(M_0,M_1>0\)，令 \(G(z)=F(z)M_0^{z-1}M_1^{-z}\)，幂以实对数定义。对 \(\varepsilon>0\)，函数 \(G(z)e^{\varepsilon(z^2-1)}\) 在两条竖边的模不超过一；在上下水平边 \(\operatorname{Im}z=\pm R\) 上，指数的模不超过 \(e^{-\varepsilon R^2}\)，故当 \(R\) 足够大时也不超过一。对矩形应用最大模原理，再令 \(R\to\infty\)、\(\varepsilon\downarrow0\)。某个 \(M_i=0\) 时，先以 \(M_i+\delta\) 替代，令 \(\delta\downarrow0\)。
\end{proof}
一个\term[jiexi-zu]{解析族}[analytic family]是复参数下的算子族，其对测试输入、测试输出的矩阵元为全纯函数。使用三线估计时必须检查边界范数和条带内增长；“公式含有 \(z\)”不是解析族的充分条件。下面先用固定算子和变化的测试函数，避免讨论一般算子的定义域变化。

\subsection{Riesz–Thorin 定理}
\label{b2:subsec:170402}

\begin{theorem}[Riesz--Thorin]\label{b2:thm:17-riesz-thorin}
设 \(T\) 是定义在两个端点 \(L^{p_0}(X)+L^{p_1}(X)\) 上的兼容线性算子，底空间为 \(\sigma\)-有限测度空间。设 \(1\leq p_i,q_i\leq\infty\)，且 \(\|Tf\|_{q_i}\leq A_i\|f\|_{p_i}\)。若
\[
 \dfrac1p=\dfrac{1-\theta}{p_0}+\dfrac\theta{p_1},\qquad
 \dfrac1q=\dfrac{1-\theta}{q_0}+\dfrac\theta{q_1},
\]
则 \(\|Tf\|_q\leq A_0^{1-\theta}A_1^\theta\|f\|_p\)。
\end{theorem}
\begin{proof}
先取有限测度支集上的简单函数 \(f,g\)，分别归一化为 \(\|f\|_p=\|g\|_{q'}=1\)，并先讨论 \(p,q'<\infty\)。在非零值处写
\[
 f_z=\dfrac f{|f|}\,|f|^{p((1-z)/p_0+z/p_1)},\quad
 g_z=\dfrac g{|g|}\,|g|^{q'((1-z)/q_0'+z/q_1')}.
\]
零值处定义为零。它们是有限个固定简单函数的解析系数组合。于是 \(F(z)=\int(Tf_z)g_z\) 全纯；在竖边 \(i=0,1\) 上有 \(\|f_{i+it}\|_{p_i}=\|g_{i+it}\|_{q_i'}=1\)，无穷指数时对应模恒为一。端点 Hölder 得到边界常数 \(A_i\)。条带内各系数的模只依赖实部，因此 \(F\) 有界。三线估计在 \(z=\theta\) 给出双线性配对的上界。

对所有简单 \(g\) 取上确界，\(L^q\) 对偶范数公式得到 \(\|Tf\|_q\) 的估计；该公式也适用于 \(q=1,\infty\)，只需用简单函数近似符号或在接近本性上确界的有限测度集合上测试。若 \(p=\infty\)，两个源端点都是无穷，直接对固定 \(f_z=f\) 作同样论证，或者用有限端点目标指数的对偶插值。一般有限 \(p\) 输入由简单函数稠密性延拓；原端点定义与延拓在公共部分相同，因而没有另造一个不兼容的算子。
\end{proof}
对解析算子族 \(T_z\)，在相同构造中把 \(T\) 换为 \(T_z\)，只要所有测试矩阵元在条带中满足同样的增长条件，结论相同。若增长至多 \(e^{c|\operatorname{Im}z|}\)，先乘 \(e^{\delta(z-\theta)^2}\) 抑制增长，得到带有限常数的估计。这一版本将在曲面限制的解析插值证明中使用。

\subsection{Hausdorff–Young 再论}
\label{b2:subsec:170403}

\begin{corollary}[Hausdorff--Young]\label{b2:thm:17-hausdorff-young}
在本卷 \(e^{-2\pi\mathrm ix\cdot\xi}\) 的约定下，\(1\leq p\leq2\) 时
\[
 \|\widehat f\|_{p'}\leq\|f\|_p.
\]
\end{corollary}
\begin{proof}
Fourier 变换从 \(L^1\) 到 \(L^\infty\) 的范数至多一；\chapref{b2:ch:04}的 Plancherel 等式给出 \(L^2\) 等距性。以这两个兼容定义应用 Riesz--Thorin。对 \(1<p\leq2\)，变换先在 \(L^1\cap L^2\cap L^p\) 上定义，再依估计延拓；不宣称任意 \(L^p\) 函数的绝对积分逐点存在。
\end{proof}
插值给出常数一，却没有证明这在中间指数最优。Gaussian 的可计算范数已经提示最优常数问题比有界性更细。逆变换具有相同的两个端点，因而同样得到 \(L^p\) 到 \(L^{p'}\) 的估计。

\subsection{插值思想}
\label{b2:subsec:170404}

两种插值处理的是不同信息。Marcinkiewicz 使用大小分解和分布函数，适用于极大函数这样的次线性算子，代价是通常不精确的常数。Riesz--Thorin 使用线性叠加及复参数，只需强型端点，常数按对数凸性组合。对奇异积分，常用弱 \((1,1)\) 和强 \((2,2)\) 先得到 \(1<p<2\)，再以伴随的相同估计补齐 \(p>2\)。

\ProseParagraphBreak

若一个估计含尺度参数 \(R\)，端点常数写成 \(A_i(R)\)，插值后的幂为 \(A_0(R)^{1-\theta}A_1(R)^\theta\)。例如两个端点分别增长为 \(R^a\) 和 \(R^b\)，中间指数的增长为 \(R^{(1-\theta)a+\theta b}\)。这项尺度核算必须在对 \(R\) 求和以前完成；一个在每个尺度上成立但常数不一致的估计，不能直接推出全算子有界。

% 扩写定位：b2:subsec:170404
为了看清“把输入分配到两个端点”与函数空间之间的联系，可以引入一个直接可计算的量。对一对嵌入同一线性空间的 Banach 空间 \(X_0,X_1\)，以及 \(f\in X_0+X_1\)，定义
\[
 K(t,f;X_0,X_1)\coloneqq
 \inf_{f=f_0+f_1}\bigl(\|f_0\|_{X_0}+t\|f_1\|_{X_1}\bigr),\qquad t>0.
\]
这个下确界允许分解随 \(t\) 改变。仅取固定分解通常丢失端点附近的信息。在 Lebesgue 空间的一对 \(L^1,L^\infty\) 中，它具有完整的重排描述。
\begin{proposition}\label{b2:prop:17-k-functional}
对 \(f\in L^1+L^\infty\)，有
\[
 K(t,f;L^1,L^\infty)=\int_0^t f^*(s)\,\mathrm ds.
\]
若线性算子 \(T:X_0+X_1\to Y_0+Y_1\) 在两端点的范数分别至多 \(A_0,A_1>0\)，则
\[
 K(t,Tf;Y_0,Y_1)\leq A_0K(tA_1/A_0,f;X_0,X_1).
\]
\end{proposition}
\begin{proof}
设 \(f=f_0+f_1\)。任意 \(|E|\leq t\) 都满足
\(\int_E|f|\leq\|f_0\|_1+t\|f_1\|_\infty\)。
在无原子的 Lebesgue 空间，取超水平集合，并在必要时取水平集合的一部分，便有
\(\sup_{|E|\leq t}\int_E|f|=\int_0^t f^*\)。当超水平集合或水平值不能恰好取到最优时，用逼近集合取极限即可。因此右侧不超过每个分解的代价。

反向令 \(a=f^*(t)\)，以保持相位的截断
\[
 f_1(x)=
 \begin{cases}f(x),&|f(x)|\leq a,\\a f(x)/|f(x)|,&|f(x)|>a,\end{cases}
 \qquad f_0=f-f_1.
\]
这里 \(f_1\) 有界，且层蛋糕给出
\(\|f_0\|_1=\int_0^\infty(f^*(s)-a)_+\,\mathrm ds\)。
由重排的单调性，这等于 \(\int_0^t f^*(s)\,\mathrm ds-ta\)，水平集合上的等号不影响积分。端点 \(a=0\) 直接同样处理；若遇到下确界的极限值，先取略大的截断高度再取极限。故该分解的代价就是所求积分。

对于算子估计，每个分解产生 \(Tf=Tf_0+Tf_1\)，代价至多 \(A_0\|f_0\|_{X_0}+tA_1\|f_1\|_{X_1}\)。对全部分解取下确界即可。
\end{proof}
由此可以定义 \(0<\theta<1\)、\(1\leq q\leq\infty\) 时的实插值量
\[
 \|f\|_{\theta,q;K}=
 \left(\int_0^\infty[t^{-\theta}K(t,f)]^q\frac{\mathrm dt}{t}\right)^{1/q},
\]
无穷指数取上确界。对 \(L^1,L^\infty\)，令 \(1/p=1-\theta\)，它与前面 \(L^{p,q}\) 的拟范数等价。下面验证这一说法，而不把它当作一个新的黑箱。

\ProseParagraphBreak

单调性给出 \(K(t,f)\geq t f^*(t)\)，所以 Lorentz 量不超过 \(K\) 量。反向置 \(h(s)=s^{1/p}f^*(s)\)，则
\[
 t^{-1/p'}K(t,f)=
 \int_0^1u^{-1/p}h(tu)\,\mathrm du.
\]
在乘法测度 \(\mathrm dt/t\) 中，变换 \(t\mapsto tu\) 保持范数；Minkowski 因而给出上界
\(\int_0^1u^{-1/p}\,\mathrm du\cdot\|h\|_{L^q(\mathrm dt/t)}
=p'\|f\|_{L^{p,q}}\)。
在 \(q=\infty\) 时直接取上确界也是同一证明。算子的 \(K\) 估计于是同时传递整个高度分布的控制。这里没有推出任意两个弱型端点间的最佳 Lorentz 指数；弱型端点不是已经给定的 Banach 范数端点，其使用还需另行检查。

\ProseParagraphBreak

对两个端点都给定强型的线性算子，代入上式并换变量 \(tA_1/A_0\)，得到插值常数 \(A_0^{1-\theta}A_1^\theta\)。这个乘法结构与 Riesz--Thorin 相同，而证明只使用分解、重排和 Minkowski。前者自然保持同一个中间 Lebesgue 指数，后者自然看到 Lorentz 的第二个指数；两种方法的适用条件和产生的信息因此各有侧重。

\ProseParagraphBreak

几种插值路线在同一个算子上可以互相核对。对可积核卷积，两个同指数端点由 Young 给出，Riesz--Thorin 重现整个强型区间；若核只在临界弱空间中，集合积分与大值分解更合适；若目标是让输入带有一个细于 \(L^p\) 的第二指数，则 \(K\) 泛函直接记录高度分布。选择方法的依据是手中实际掌握的端点，而不是插值定理名字的先后次序。

\ProseParagraphBreak

实际证明还应先检查共同定义域。比如 Fourier 变换同时有 \(L^1\to L^\infty\) 与 \(L^2\to L^2\) 的界，这两种对象在 \(L^1\cap L^2\) 上一致，才可以把它看成同一个待插值算子。若两端的公式只分别定义在互不相容的空间上，即使形式上写出了两个范数界，也没有一个可供插值的共同线性算子。对奇异积分，这个相容性通常由测试函数上的主值公式与平方延拓来建立。

\ProseParagraphBreak

输出是否线性同样重要。极大算子把所有尺度的绝对平均取上确界，一般不是线性的；把输入换成复幂函数后，不能期待极大函数的矩阵元对参数解析。因此本章对它采用分布函数插值。相比之下，固定的 Fourier 变换或固定卷积是线性的，解析变化发生在测试输入与测试输出中，配对仍可验证为全纯函数。假如把“绝对值”“上确界”和“积分”在解析证明中任意交换，便可能把非解析量误放入三线引理。

\ProseParagraphBreak

常数也可以作为一个有效检查。若用两端范数 \(A_0,A_1\) 推导中间界，缩放算子 \(T\mapsto cT\) 后，中间界必须线性乘以 \(|c|\)。Riesz--Thorin 的 \(A_0^{1-\theta}A_1^\theta\) 正好具有这种齐次性；Marcinkiewicz 的未优化分割可能给出较粗的常数，但仍应在适当重选高度阈值后保持齐次。若某个推导出现两个端点常数相乘、总指数大于一，通常意味着阈值没有随算子大小调整，或只是用了一个合法但很粗的非齐次上界。

\ProseParagraphBreak

插值不能恢复被端点假设遗忘的信息。两个范数界可能来自完全不同的机制，插值只利用其共同的线性或次线性结构；它不会证明核的抵消、主值存在性、最优常数或极值函数分类。例如 Hausdorff--Young 的常数一来自基本插值，而 Gaussian 测试揭示中间指数还有更细的最佳常数问题。正向上界与某类测试函数给出的下界应分别记录，只有两者相等才完成最优常数的计算。

\ProseParagraphBreak

最后，开放指数区间中的有界性并不总能通过取极限传到端点。若估计常数在 \(p\to1\) 时发散，对每个固定输入得到的范数序列就没有统一控制；即使常数不发散，也要说明输入空间与输出空间随指数变化时采用了什么收敛。极大函数的具体反例已经给出最直接的检验：它在每个 \(p>1\) 上有界，然而一个单位区间的像有不可积尾部。弱端点的提出正是为了保留这种有限而有意义的控制。
